\documentclass[11pt,a4paper]{article}
\usepackage[margin=2.4cm]{geometry}
\usepackage{amsmath,amssymb,booktabs,hyperref}
\newcommand{\dt}{\Delta t}
\newcommand{\avg}[1]{\{\!\{#1\}\!\}}
\newcommand{\jump}[1]{[\![#1]\!]}
\newcommand{\code}[1]{\texttt{#1}}
\title{A frozen split of the DGSEM Euler equations\\with the pressure gradient in product form}
\author{MultiRate notes}
\date{\today}
\begin{document}
\maketitle

\begin{abstract}
For split-explicit (MIS) time integration with a \emph{frozen} right-hand side, $N(q)=F(q,q)$, the
first argument $y$ is frozen at the stage values and the second $z$ evolves in the fast loop.
We describe a DGSEM discretization of the compressible Euler equations in the variables
$(\rho,\ m=\rho\mathbf u,\ \Theta=\rho\theta)$ in which this split is natural: the pressure gradient
is written in product form $\nabla p = P'(\Theta)\nabla\Theta$, the potential temperature $\theta$ and
the derivative $P'=\mathrm dp/\mathrm d\Theta$ are frozen exactly where they appear, and the interface
flux is a LMARS variant written directly in the momentum $m$. The fast part is then \emph{linear} in
$z$ with coefficients frozen at $y$, and $F(q,q)=N(q)$ holds by construction.
\end{abstract}

\section{Equations}
With $\Theta=\rho\theta$ the pressure is a function of $\Theta$ alone,
\[
p=P(\Theta)=p_0\Big(\frac{R\,\Theta}{p_0}\Big)^{\gamma},\qquad
P'(\Theta)=\frac{\gamma p}{\Theta}=\frac{c^2}{\theta},\qquad c^2=\frac{\gamma p}{\rho}.
\]
The Euler equations with gravity, the pressure gradient in product form:
\begin{align}
\partial_t\rho+\nabla\cdot m&=0,\label{eq:rho}\\
\partial_t\Theta+\nabla\cdot(\theta\,m)&=0,\label{eq:theta}\\
\partial_t m+\nabla\cdot(m\otimes\mathbf u)+P'(\Theta)\,\nabla\Theta+\rho g\,\mathbf e_z&=0.\label{eq:mom}
\end{align}
Mass and $\Theta$ are in conservation form; the momentum equation is not (the pressure term is a
non-conservative product), which is of no concern with gravity present.
The equations are discretized as they stand, without a reference state and without subtracting a
hydrostatic background (Section~\ref{sec:balance}).

\section{The frozen form}
Freeze $\theta$ and $P'$ at $y$, everything else from $z$:
\begin{align*}
F_\rho(y,z)&=-\nabla\cdot m_z,\\
F_\Theta(y,z)&=-\nabla\cdot(\theta_y\,m_z),\\
F_m(y,z)&=\underbrace{-\nabla\cdot(m_y\otimes\mathbf u_y)}_{\text{slow }S(y)}
\;\underbrace{-\,P'(\Theta_y)\nabla\Theta_z-\rho_z g\,\mathbf e_z}_{\text{fast}}.
\end{align*}
$F(q,q)$ is \eqref{eq:rho}--\eqref{eq:mom}. The fast part
\[
F(y,z)-S(y)=L(y)\,z
\]
is linear in $z=(\rho,m,\Theta)_z$; its operator $L(y)$ has the coefficients $\theta_y$ and $P'(\Theta_y)$
only. Nothing else has to be frozen: the momentum enters linearly, the pressure only through $\nabla\Theta$.
A stage of MIS integrates
\[
\frac{\mathrm dZ}{\mathrm d\tau}=\sum_j\frac{\beta_{ij}}{d_i}\big(S(Y_j)+L(Y_j)Z\big)+\sum_j\frac{\gamma_{ij}}{d_i\dt}(Y_j-y^n),
\]
i.e.\ with the weighted coefficients $\bar\theta=\sum_j b_{ij}\theta(Y_j)$, $\bar P'=\sum_j b_{ij}P'(\Theta(Y_j))$,
$b_{ij}=\beta_{ij}/d_i$ (since $L$ is linear in the coefficients), one linear operator per stage.

\section{DGSEM discretization}

\subsection{Notation}
Quadrilateral elements mapped from $[-1,1]^2$, $(\xi,\eta)\mapsto\mathbf x$, Jacobian $J$, metric terms
$J\xi_x, J\xi_z, J\eta_x, J\eta_z$ (discretely satisfying the metric identities), tensor-product LGL nodes of
degree $p$, collocated quadrature with weights $w_i$, differentiation matrix $D$ (SBP:
$WD+D^TW=B=\mathrm{diag}(-1,0,\dots,0,1)$). Contravariant components of a vector $v$:
$\tilde v^\xi=J\xi_x v_x+J\xi_z v_z$, $\tilde v^\eta=J\eta_x v_x+J\eta_z v_z$. At a face node with outward
unit normal $\mathbf n$ and surface element $s$ (so that $s\,\mathbf n$ is the contravariant normal),
$v_n=\mathbf v\cdot\mathbf n$; $L$/$R$ denote the inner and the outer trace, $\avg a=\tfrac12(a_L+a_R)$,
$\jump a=a_R-a_L$.

\subsection{Volume terms}
In flux-differencing form along the $\xi$ lines (the $\eta$ lines likewise), node $a=(i,j)$, partner
$b=(k,j)$, with the two-point contravariant normal $\tilde n_{ab}=\tfrac12\big((J\xi_x)_a+(J\xi_x)_b,\ (J\xi_z)_a+(J\xi_z)_b\big)$:
\begin{align*}
(\text{mass})_a&=\frac1{J_a}\sum_k 2D_{ik}\ \avg{m_z}_{ab}\cdot\tilde n_{ab},\\
(\Theta)_a&=\frac1{J_a}\sum_k 2D_{ik}\ \avg{\theta_y}_{ab}\,\avg{m_z}_{ab}\cdot\tilde n_{ab},\\
(\text{advection, slow})_a&=\frac1{J_a}\sum_k 2D_{ik}\ \big(\avg{m_y}_{ab}\cdot\tilde n_{ab}\big)\avg{\mathbf u_y}_{ab},\\
(\text{pressure})_a&=P'(\Theta_{y,a})\ \frac1{J_a}\sum_k D_{ik}\,\Theta_{z,b}\,(J\xi_x,J\xi_z)_b
\quad\big(\text{the conservative metric form of }\nabla\Theta\big),\\
(\text{gravity})_a&=\rho_{z,a}\,g\,\mathbf e_z .
\end{align*}
Since $\sum_kD_{ik}=0$ the mass term is the strong-form divergence of $m_z$; for the $\Theta$ term the
two-point mean $\avg{\theta_y}$ is the only frozen coefficient. The pressure term is the nodal product of
the frozen $P'_a$ with the DG gradient of $\Theta_z$ (its surface part in the next section).

\subsection{Interface flux: LMARS in the momentum}
At a face the normal acoustics of \eqref{eq:theta}--\eqref{eq:mom}, frozen at a face state
$(\theta_f, c_f)$, is the linear system
\[
\partial_t\Theta+\theta_f\,\partial_n m_n=0,\qquad
\partial_t m_n+\frac{c_f^2}{\theta_f}\,\partial_n\Theta=0,
\]
with characteristic speeds $\pm c_f$ and the Riemann invariants $m_n\pm (c_f/\theta_f)\Theta$. Its exact
Riemann solution at the face is
\begin{equation}
m^*=\avg{m_n}-\frac{c_f}{2\theta_f}\,\jump{\Theta},\qquad
\Theta^*=\avg{\Theta}-\frac{\theta_f}{2c_f}\,\jump{m_n}.\label{eq:riemann}
\end{equation}
This is LMARS written in $m$ instead of $\mathbf u$: with $\jump p\approx P'\jump\Theta=(c^2/\theta)\jump\Theta$,
$m^*=\avg{m_n}-\jump p/(2c)$ and $p^*=P'\Theta^*=\avg p-\tfrac12 c\,\jump{m_n}$, i.e.\ the LMARS
interface velocity and pressure $w^*=\avg{u_n}-\jump p/(2\rho c)$, $p^*=\avg p-\tfrac12\rho c\jump{u_n}$ with the
density factor absorbed into the momentum. No density upwinding $\rho_{\rm up}$ and no $1/\rho$ appear:
$m^*$ is linear in $(m,\Theta)_z$.

The interface fluxes (per unit contravariant normal, outward from $L$):
\begin{align*}
\hat f_\rho&=m^*,\qquad \hat f_\Theta=\theta_{\rm up}\,m^*,\qquad
\hat f_m^{\rm adv}=m^*_y\,\mathbf u_{{\rm up},y}\quad(\text{slow}),
\end{align*}
$\theta_{\rm up}$, $\mathbf u_{\rm up}$ the values of the upwind side, the side chosen by the sign of $m^*(y)$.
In strong form the surface term at a face node of element $L$ is
\[
\frac{s}{J\,w}\Big(\hat f-f_L\cdot\mathbf n\Big),\qquad f_{\rho,L}\cdot\mathbf n=m_{n,L},\quad f_{\Theta,L}\cdot\mathbf n=\theta_L m_{n,L},
\]
and for the non-conservative pressure term the lifted jump to the interface value,
\[
\frac{s}{J\,w}\,P'(\Theta_{y,L})\,\big(\Theta^*-\Theta_{z,L}\big)\,\mathbf n .
\]
\paragraph{Frozen face coefficients.} $\theta_f=\avg{\theta_y}$, $c_f=\avg{c_y}$ (or $\max$), the upwind side and
$\theta_{\rm up}$ from $y$. With $z$ evolving, $m^*$ and $\Theta^*$ are linear in $z$; the face part of the fast
operator needs $c_f/(2\theta_f)$, $\theta_f/(2c_f)$, $\theta_{\rm up}$ and the nodal $P'_L$, the slow part
$m^*(y)\,\mathbf u_{\rm up}(y)$.
\paragraph{Contact discontinuities.} Because $p=P(\Theta)$, a jump in $\Theta$ is a jump in $p$. A contact
(constant $p$, jump in $\theta$ and $\rho$) has $\jump\Theta=0$ and gets no acoustic dissipation in $m^*$;
$\theta$ is upwinded in $\hat f_\Theta$, $\rho$ is carried by the central $\avg{m_n}$.
\paragraph{Walls.} Mirror state $m_{n,R}=-m_{n,L}$, $\Theta_R=\Theta_L$: $m^*=0$ and
$\Theta^*=\Theta_L+\theta_f m_{n,L}/c_f$, so the pressure surface term is
$P'_L(\Theta^*-\Theta_L)=c\,m_{n,L}$ (with $P'\theta=c^2$), the LMARS wall pressure $p+\rho c\,u_n$.

\subsection{Hydrostatic balance}\label{sec:balance}
No background state is subtracted. A hydrostatic state given at the nodes, continuous across faces,
has no jumps ($\Theta^*=\Theta$, $m^*=0$ at interior faces and walls), so the discrete momentum residual
is the volume term alone, $r_a=P'(\Theta_a)\,(\mathcal G\Theta)_a+\rho_a g\,\mathbf e_z$, i.e.\ the
pointwise error of the degree-$p$ differentiation of $\Theta(z)$, of order $\Delta z^{p}$ per element
(spectral in $p$). An initial state in exact discrete balance is obtained by solving
$P'(\Theta_a)(\mathcal G\Theta)_a=-\rho_a g$ for $\rho_a$ at given $\Theta$ (nodewise, since the operator
acts on $\Theta$ only), which needs no change of the scheme. Section~\ref{sec:wb} gives a variant that
is exactly balanced for the isothermal and the constant-$\theta$ atmosphere, still without subtraction.

\section{Well-balanced variant}\label{sec:wb}
Artiano, Knoth, Spichtinger, Ranocha (\emph{Structure-preserving high-order methods for the
compressible Euler equations in potential temperature formulation for atmospheric flows},
arXiv:2509.10311) obtain well-balanced DGSEM schemes without any background state: gravity
$-\rho\nabla\phi$ ($\phi=gz$ the geopotential, a given function) is discretized as a
non-conservative two-point term $\bar\rho_{ab}\jump{\phi}_{ab}$, the pressure by the two-point flux
$\avg p$, and the mean $\bar\rho$ is chosen such that $\bar\rho\jump\phi=-\jump p$ holds exactly for the
hydrostatic state:
\[
\text{isothermal: }\ \bar\rho=\avg\rho_{\log}=\frac{\jump\rho}{\jump{\log\rho}},\qquad
\text{constant }\theta:\ \ \bar\rho=\avg\rho_\gamma=\frac{\gamma-1}{\gamma}\,
\frac{\jump{\rho^\gamma}}{\jump{\rho^{\gamma-1}}}.
\]
(Isothermal: $p=RT\rho$, $\jump{\log\rho}=-\jump\phi/(RT)$, so $\avg\rho_{\log}\jump\phi=-RT\jump\rho=-\jump p$.
Constant $\theta$: $p=K\theta^\gamma\rho^\gamma$ and $\tfrac{\gamma}{\gamma-1}K\theta^\gamma\rho^{\gamma-1}+\phi=\mathrm{const}$,
so $\avg\rho_\gamma\jump\phi=-K\theta^\gamma\jump{\rho^\gamma}=-\jump p$.) Since $\phi$ is continuous, gravity has no
surface term.

\paragraph{Product form with two-point coefficients.} The product form fits this exactly if $\mathrm dp/\mathrm d\Theta$
is frozen as a \emph{two-point secant} where it appears. Along a $\xi$ line (pair $a=(i,j)$, $b=(k,j)$,
$\avg{J\xi}_{ab}=\tfrac12\big((J\xi_x,J\xi_z)_a+(J\xi_x,J\xi_z)_b\big)$; $\eta$ lines likewise):
\begin{align*}
(\text{pressure})_a&=\frac1{J_a}\sum_k D_{ik}\,P'_{ab}\,\jump{\Theta}_{ab}\,\avg{J\xi}_{ab},\qquad
P'_{ab}=\frac{\jump p_{ab}}{\jump\Theta_{ab}}\quad\big(\to P'(\Theta)\text{ for }\Theta_a=\Theta_b\big),\\
(\text{gravity})_a&=\frac1{J_a}\sum_k D_{ik}\,\bar\rho_{ab}\,\jump{\phi}_{ab}\,\avg{J\xi}_{ab}.
\end{align*}
With the secant, $P'_{ab}\jump\Theta_{ab}=\jump p_{ab}$, and since $2\avg p_{ab}=2p_a+\jump p_{ab}$ and
$\sum_{\xi,\eta}\sum_kD_{ik}\avg{J\xi}_{ab}=0$ (discrete metric identities), the pressure term is identical to the
conservative flux differencing $\frac1{J_a}\sum_k2D_{ik}\avg p_{ab}\avg{J\xi}_{ab}$: the product form loses nothing.
For the hydrostatic state every pair vanishes separately, $P'_{ab}\jump\Theta_{ab}+\bar\rho_{ab}\jump\phi_{ab}
=\jump p_{ab}+\bar\rho_{ab}\jump\phi_{ab}=0$, so the volume residual is zero for any degree and mesh. At the faces
the state is continuous and at rest: $\jump\Theta=0$, $m^*=0$, $\Theta^*=\Theta_L$, and the pressure surface term
vanishes as well. With the face secant $P'_{L*}=\big(P(\Theta^*)-P(\Theta_L)\big)/(\Theta^*-\Theta_L)$ the surface term
$P'_{L*}(\Theta^*-\Theta_L)=p(\Theta^*)-p_L$ is the conservative one, too.

\paragraph{Frozen form.} Freeze the coefficients where they appear, from $y$:
\[
P'_{ab}(y)=\frac{\jump{p_y}_{ab}}{\jump{\Theta_y}_{ab}},\qquad
R_{ab}(y)=\frac{\bar\rho_{ab}(y)}{\avg{\rho_y}_{ab}},\qquad P'_{L*}(y),
\]
and take $\jump{\Theta_z}_{ab}$, $\Theta^*_z-\Theta_{z,L}$ and $R_{ab}\avg{\rho_z}_{ab}\jump\phi_{ab}$ live. The fast part stays
linear in $z$, $F(q,q)=N(q)$ holds, and at $y=z=q$ hydrostatic the scheme is exactly balanced. This is the
same secant as in the FV frozen split (\code{FV\_FROZEN\_PRESSURE = :derivative}); the frozen gravity ratio $R$
is the one of the present conservative implementation.

\paragraph{Other stratifications.} For a general hydrostatic profile (e.g.\ constant Brunt--V\"ais\"al\"a
frequency, the gravity-wave tests) neither mean is exact; the residual is then of the truncation
order of the mean, $\bar\rho-\rho^{\rm exact}_{ab}=O(\Delta z^2)$ per pair. Choosing $\avg\rho_\gamma$ (exact for
constant $\theta$) is the natural default for the troposphere; a profile-adapted mean would have to be a
mean of $\rho$ alone (horizontal pairs have $\jump\phi=0$ and must not see it).

\section{Summary of the split}
\begin{center}
\footnotesize
\begin{tabular}{lll}
\toprule
term & frozen from $y$ & live from $z$\\
\midrule
volume mass flux & -- & $\avg{m}$\\
volume $\Theta$ flux & $\avg{\theta}$ & $\avg{m}$\\
pressure (volume and surface) & nodal $P'=\gamma p/\Theta$ (well-balanced: secants $P'_{ab}$, $P'_{L*}$) & $\mathcal G\Theta$ ($\jump\Theta_{ab}$), $\Theta^*$\\
gravity & -- (well-balanced: $R_{ab}=\bar\rho_{ab}/\avg{\rho}_{ab}$) & $\rho$ ($\avg\rho_{ab}$)\\
face $m^*$, $\Theta^*$ & $c_f/\theta_f$, $\theta_f/c_f$ & $\avg{m_n},\jump\Theta,\avg\Theta,\jump{m_n}$\\
face $\Theta$ flux & $\theta_{\rm up}$, upwind side & $m^*$\\
\midrule
momentum advection (slow) & everything & --\\
\bottomrule
\end{tabular}
\end{center}

\section{Properties}
\begin{enumerate}
\item $F(q,q)=N(q)$ by construction: $N$ \emph{is} the product-form DGSEM with these fluxes.
\item The fast part is linear in $z$. A HEVI fast loop needs one linear solve per column with the frozen
  vertical coefficients (no Newton iteration for $p(\Theta)$, unlike the conservative form, where the
  nodal $p(\Theta_z)$ is nonlinear). The coefficients of a stage are the $b_{ij}$-weighted means of
  $\theta$, $P'$, $c_f/\theta_f$ over the stage values.
\item Mass and $\Theta$ are conserved; momentum is not (non-conservative pressure).
\item Linear acoustics: for constant frozen coefficients the volume part is skew-adjoint in the energy
  $\tfrac12\sum w J\big(|m|^2+(c^2/\theta^2)\Theta^2\big)$ (SBP), and the exact Riemann solver
  \eqref{eq:riemann} adds the dissipation
  $-\sum_{\rm faces} s\,w\,\tfrac{c}{2}\big(\jump{m_n}^2+\tfrac{c^2}{\theta^2}\jump\Theta^2\big)\le0$
  (upwind flux of the symmetrized variables $(m_n,\,c\,\Theta/\theta)$).
  Variable coefficients add terms of the size of their variation, as usual.
\item Only $\theta$ and $P'$ (and the face sound speed) are frozen, the ``coefficients'' of the
  fast operator; no part of the stiff response sits in the frozen argument. This is the rule of thumb
  found for the FV frozen split (freezing $p/\Theta$ instead of $\mathrm dp/\mathrm d\Theta$ was unstable).
\end{enumerate}

\section{Relation to the current implementation and open points}
\begin{itemize}
\item \code{src/DGSEM/Euler2D/frozen.jl} splits the \emph{conservative} DGSEM (pressure $\{p\}$ in the
  two-point momentum flux, gravity with a logarithmic mean of $\rho$, LMARS in the velocity): there
  the frozen coefficients are $A/\rho_a$, $A/\rho_b$, $\rho_{\rm up}/\rho$, the ratio of the logarithmic
  to the arithmetic mean, and the fast part keeps the nonlinear nodal $p(\Theta_z)$. The product form
  removes all $1/\rho$ factors and makes the fast part linear.
\item MIS stability. In both forms the $\Theta$ flux is $\theta_y m_z$: linearized about a wind $\bar u$,
  its fast part is $\bar\theta\,\delta m=\bar\theta(\bar u\,\delta\rho+\bar\rho\,\delta u)$ and its slow part
  $\bar m\,\delta\theta$. Density is advected in the fast part, $\Theta$ advection is split. On the
  column Fourier blocks this coupling limits the stability of MIS methods
  (\code{MethodDesign/app\_stability.jl}; methods designed for it: \code{MIS\_DG53}). The analysis has to
  be repeated for the product form.
\item To do: implementation (\code{frozen\_ncoef}, \code{frozen\_data!}, \code{frozen\_fast!} with the
  product-form $N$), test $F(q,q)=N(q)$, the frozen HEVI column solve, gravity-wave and terrain tests,
  and a comparison of the product form with the conservative DGSEM.
\end{itemize}
\end{document}
